%! Equivalent Representations of Trigonometric Functions — Unit 3, Lesson 3.12 — Slides (Beamer)
\documentclass[aspectratio=169, 11pt]{beamer}
\usepackage{precalculus-beamer}   % provides \navyheader{} and \sectionlabel[color]{}

\begin{document}

% TITLE SLIDE
{
\setbeamercolor{background canvas}{bg=navy}
\begin{frame}[plain]
  \vspace{2.8cm}\hspace{0.55cm}%
  \begin{minipage}{0.9\textwidth}
    {\color{goldacc}\bfseries\small LESSON 3.12}\\[0.5cm]
    {\color{white}\bfseries\LARGE Equivalent Representations of Trigonometric Functions}\\[1.2cm]
    {\color{skymid}\small AP Precalculus~$\cdot$~Unit 3}
  \end{minipage}
\end{frame}
}

%----------------------------------------------------------------------
% Frame 1 — Pythagorean Identity
%----------------------------------------------------------------------
\begin{frame}[t, plain]
  \navyheader{LO 3.12.A --- The Pythagorean Identity}

  \begin{columns}[T]
  \column{0.48\textwidth}
    \begin{block}{Derivation from the Unit Circle}
      A point on the unit circle: $(\cos\theta,\, \sin\theta)$\\[4pt]
      Distance from origin $= 1$ (the radius)\\[4pt]
      By the Pythagorean Theorem:
      $$\cos^2\theta + \sin^2\theta = 1$$
    \end{block}

    \vspace{6pt}
    \begin{block}{Three Forms}
      \begin{itemize}
        \item $\sin^2\theta + \cos^2\theta = 1$
        \item Divide by $\cos^2\theta$: \quad $1 + \tan^2\theta = \sec^2\theta$
        \item Divide by $\sin^2\theta$: \quad $\cot^2\theta + 1 = \csc^2\theta$
      \end{itemize}
    \end{block}

  \column{0.48\textwidth}
    \begin{block}{Key Uses}
      \begin{itemize}
        \item Substitute to \textbf{reduce} two trig functions to one
        \item Find a missing trig value given one value and the quadrant
        \item Simplify complex trig expressions
      \end{itemize}
    \end{block}

    \vspace{6pt}
    \textbf{Example:} If $\cos\theta = \tfrac{3}{5}$ (Q IV), find $\sin\theta$.

    \vspace{4pt}
    $\sin^2\theta = 1 - \left(\tfrac{3}{5}\right)^2 = \tfrac{16}{25}$

    $\Rightarrow \sin\theta = -\tfrac{4}{5}$ (negative in Q IV)
  \end{columns}
\end{frame}

%----------------------------------------------------------------------
% Frame 2 — Sum and Difference Identities
%----------------------------------------------------------------------
\begin{frame}[t, plain]
  \navyheader{LO 3.12.B --- Sum and Difference Identities}

  \begin{columns}[T]
  \column{0.50\textwidth}
    \begin{block}{The Four Identities}
      \begin{itemize}
        \item $\sin(\alpha+\beta) = \sin\alpha\cos\beta + \cos\alpha\sin\beta$
        \item $\sin(\alpha-\beta) = \sin\alpha\cos\beta - \cos\alpha\sin\beta$
        \item $\cos(\alpha+\beta) = \cos\alpha\cos\beta - \sin\alpha\sin\beta$
        \item $\cos(\alpha-\beta) = \cos\alpha\cos\beta + \sin\alpha\sin\beta$
      \end{itemize}
    \end{block}

    \vspace{4pt}
    \textbf{Memory tip for cosine:} The sign \emph{flips} ---
    sum identity has $-$, difference has $+$.

  \column{0.46\textwidth}
    \begin{block}{Worked Example --- $\sin 75^\circ$}
      Write $75^\circ = 45^\circ + 30^\circ$:
      \begin{align*}
        &\sin 45^\circ\cos 30^\circ + \cos 45^\circ\sin 30^\circ \\
        &= \tfrac{\sqrt{2}}{2}\cdot\tfrac{\sqrt{3}}{2} + \tfrac{\sqrt{2}}{2}\cdot\tfrac{1}{2} \\
        &= \tfrac{\sqrt{6}}{4} + \tfrac{\sqrt{2}}{4}
         = \dfrac{\sqrt{6}+\sqrt{2}}{4}
      \end{align*}
    \end{block}

    \vspace{4pt}
    \textbf{Key idea:} We can find exact values for angles like
    $15^\circ$, $75^\circ$, $\tfrac{\pi}{12}$, $\tfrac{5\pi}{12}$
    by splitting them into standard unit-circle angles.
  \end{columns}
\end{frame}

%----------------------------------------------------------------------
% Frame 3 — Double-Angle Identities
%----------------------------------------------------------------------
\begin{frame}[t, plain]
  \navyheader{LO 3.12.B --- Double-Angle Identities}

  \begin{columns}[T]
  \column{0.48\textwidth}
    \begin{block}{Derived from Sum Identities ($\alpha = \beta = \theta$)}
      \begin{itemize}
        \item $\sin 2\theta = 2\sin\theta\cos\theta$
        \item $\cos 2\theta = \cos^2\theta - \sin^2\theta$
        \item $\cos 2\theta = 2\cos^2\theta - 1$
        \item $\cos 2\theta = 1 - 2\sin^2\theta$
      \end{itemize}
    \end{block}

    \vspace{4pt}
    \begin{block}{Useful Rearrangements}
      $$\sin^2\theta = \frac{1 - \cos 2\theta}{2}$$
      $$\cos^2\theta = \frac{1 + \cos 2\theta}{2}$$
    \end{block}

  \column{0.48\textwidth}
    \begin{block}{Why Three Cosine Forms?}
      \begin{itemize}
        \item $\cos^2\theta - \sin^2\theta$: compact symmetric form
        \item $2\cos^2\theta - 1$: eliminates $\sin\theta$
        \item $1 - 2\sin^2\theta$: eliminates $\cos\theta$
      \end{itemize}
      Choose the form that converts the equation to a \emph{single} function.
    \end{block}

    \vspace{4pt}
    \textbf{Example:} $\cos^2\theta - \sin^2\theta = 0$

    $\Rightarrow \cos 2\theta = 0$

    $\Rightarrow 2\theta = \tfrac{\pi}{2}, \tfrac{3\pi}{2}, \ldots$

    $\Rightarrow \theta = \tfrac{\pi}{4}, \tfrac{3\pi}{4}, \tfrac{5\pi}{4}, \tfrac{7\pi}{4}$
  \end{columns}
\end{frame}

%----------------------------------------------------------------------
% Frame 4 — Solving with Identities
%----------------------------------------------------------------------
\begin{frame}[t, plain]
  \navyheader{LO 3.12.C --- Solving Equations with Identities}

  \begin{columns}[T]
  \column{0.48\textwidth}
    \begin{block}{Four-Step Strategy}
      \begin{enumerate}
        \item Identify all trig functions present
        \item Choose an identity to reduce to one function
        \item Solve the resulting algebraic equation
        \item Find \textbf{all} solutions in the given interval (check both quadrants)
      \end{enumerate}
    \end{block}

    \vspace{4pt}
    \textbf{Common substitution:}

    If the equation has $\sin^2\theta$, substitute $1 - \cos^2\theta$.

    If it has $\cos^2\theta$, substitute $1 - \sin^2\theta$.

  \column{0.48\textwidth}
    \begin{block}{Worked Example}
      Solve $2\sin^2\theta + \cos\theta - 1 = 0$ on $[0, 2\pi)$

      \vspace{4pt}
      Substitute $\sin^2\theta = 1 - \cos^2\theta$:

      $2(1 - \cos^2\theta) + \cos\theta - 1 = 0$

      $-2\cos^2\theta + \cos\theta + 1 = 0$

      $2\cos^2\theta - \cos\theta - 1 = 0$

      $(2\cos\theta + 1)(\cos\theta - 1) = 0$

      $\cos\theta = -\tfrac{1}{2}$ \; or \; $\cos\theta = 1$

      \vspace{2pt}
      $\theta \in \left\{0,\; \tfrac{2\pi}{3},\; \tfrac{4\pi}{3}\right\}$
    \end{block}
  \end{columns}
\end{frame}

\end{document}
